\section{Lorentz bridge: effective kinematics for inside observers}
\label{sec:lorentz}

This bridge imports the linearized branch speeds $c_L^2 = k_s a^2/m$ and
$c_T^2 = (1-\alpha)\,k_s a^2/m$ from Paper~I (\Cref{eq:wave-speeds-derived})
and asks whether inside observers built from the same substrate material would
experience these as a universal, direction-independent Lorentz invariant kinematics.
The ambient lattice is cubic and is \emph{not} rotationally symmetric; the bridge
argument must therefore close the gap between lab-frame anisotropy and
observer-frame isotropy.

\subsection{Linear-wave cone and effective Lorentz kinematics}

The substrate carries Lorentzian sign structure only at the level of the brane action;
the ambient $\R^4$ in which the brane is embedded is symmetric Euclidean.
A Lorentz metric on excitations is therefore not postulated --- it is an emergent property
of the linear wave regime as seen by an inside observer.

From Paper~I, linearizing the equations of motion about a homogeneous background yields the
long-wavelength dispersion
\begin{equation}
  \omega_b^2(\veck)\;\approx\; c_b^2\,|\veck|^2 \quad \text{for } |\veck|\to 0,
  \label{eq:linear-dispersion}
\end{equation}
which defines an effective wave cone for each polarization branch $b$.
In regimes where a single branch dominates transport and dispersion is negligible, the wave
equation is approximately invariant under the standard Lorentz group at speed $c \equiv c_T$,
with boost parameter
\begin{equation}
  \gamma = \bigl(1-v^2/c^2\bigr)^{-1/2},
  \label{eq:lorentz-boost}
\end{equation}
as an effective symmetry of the linear propagation sector, not as a microscopic postulate.

\subsection{Dual-observer argument: lab anisotropy vs internal isotropy}
\label{para:dual-observer}

On the 6-neighbor axial-only lattice the long-wavelength wave speed is direction-dependent at
the lab-frame level --- for example $c_L([100])/c_L([111])\approx 1.075$ at the default
operating point $\alpha=0.2$.
The relativistic effective description \eqref{eq:lorentz-boost} is therefore \emph{not} a
statement about the lab frame; it is a conjecture about \emph{internal} observers, whose rods,
clocks, and signal cones are substrate-emergent.

The closing mechanism is as follows.
A rod built from substrate material is itself a strain pattern propagating at the same branch
speeds as the signal.
An inside observer's unit-length standard rod in direction $\hat{n}$ therefore has a physical
length (in the lab metric) that scales with $c(\hat{n})$ --- exactly the same direction
dependence as the signal cone in that direction.
When both the signal and the measuring rod renormalize with the same anisotropy factor, the
observed ratio (distance/signal-transit-time) is direction-independent to leading order.

This is not a tunable coincidence.
It is a structural consequence of the principle that observers and signals are made of the same
substrate material: the stiffness tensor that sets the branch speed also sets the elastic
response of a rod-shaped excitation, so anisotropy factors cancel in any local observer
measurement.

\paragraph{Falsifiable residue.}
The cancellation is exact only in the strict long-wavelength, single-branch, weak-amplitude
limit.
At finite wavelength and nonzero mixing between branches, a residual direction dependence
survives.
Its magnitude at scale $ka$ is a calculable quantity that depends on the stiffness tensor and
neighbor stencil.
\emph{Residual birefringence} --- a direction-dependent splitting of effective light speed ---
is therefore the primary falsifiable signature of the bridge: if two solitons in different
orientations measure the same effective signal speed to better than the residual predicted by the
stiffness tensor, the bridge is supported; if they do not, it fails.

\subsection{Analogue-metric viewpoint}

An equivalent viewpoint, standard in analogue-gravity systems, is that linear excitations
experience an effective spacetime interval
\begin{equation}
  ds^2 = -c^2\,dt^2 + g^{(\text{eff})}_{ij}(x,t)\,dx^i\,dx^j,
  \label{eq:effective-metric-interval}
\end{equation}
where $g^{(\text{eff})}_{ij}$ depends on the background configuration about which one
linearizes.
Two geometric objects contribute to this background:
the induced brane metric $g_{ij}(\vecX)$ from Paper~I \eqref{eq:induced-metric}, and
the branch-dependent effective coefficients of the linearized wave operator.
Slowly varying transverse profiles $u = X^4$ modify intrinsic distances at the level of the
induced metric through the term $\partial_i u\,\partial_j u$, and can therefore be read as a
background that focuses or defocuses wave propagation in close analogy to a weak gravitational
field.
Both encode how background deformations bend wave transport, as in analogue-gravity
constructions \citep{Barcelo2011AnalogueGravity,Volovik2003HeliumDroplet}.

\subsection{Regime of validity and expected deviations}
\label{subsec:lorentz-validity}

Effective Lorentz symmetry is expected to break when:
(i) multiple branches contribute comparably (mode mixing);
(ii) dispersion becomes significant at shorter wavelengths;
(iii) geometric or material nonlinearity becomes important at larger amplitudes.

In the present theory these deviations are structural consequences of a substrate whose timelike
direction is selected by the brane action rather than imposed as a metric primitive.
They delimit the domain where a relativistic effective description applies, and each regime
transition is a calculable correction to the dual-observer argument, not an unfalsifiable
escape hatch.

\subsection{Lorentz priority open problems}
\label{subsec:lorentz-open}

\begin{enumerate}
  \item \textbf{Observer-universality calculation.}
  Compute the direction-dependent wave speed on the 6-neighbor cubic stencil at finite $ka$
  and determine the magnitude of residual birefringence as a function of wavelength and $\alpha$.
  This bounds from above the scale at which effective Lorentz invariance is expected.

  \item \textbf{Two-soliton isotropy test.}
  Run two solitons in orthogonal directions and compare their effective internal metric
  (signal-transit-time ratios).
  Agreement to within the predicted residual birefringence is the decisive numerical check.

  \item \textbf{Multi-stencil tuning.}
  Explore whether shell weights on the cubic stencil can be adjusted to suppress residual
  anisotropy further, and at what cost to other bridge properties (e.g.\ wave speed splitting).
\end{enumerate}